\documentclass[a4paper,11pt]{article}
\usepackage[utf8]{inputenc}
\usepackage[T2A,T1]{fontenc}
\usepackage[bulgarian,english]{babel}
\usepackage[top=2.2cm, bottom=2.2cm, left=2.4cm, right=2.4cm]{geometry}
\usepackage{fontawesome5}
\usepackage{xcolor}
\usepackage{titlesec}
\usepackage{enumitem}
\usepackage{multicol}
\usepackage{mdframed}
\usepackage{tcolorbox}
\usepackage{tabularx}
\usepackage{array}
\usepackage{fancyhdr}
\usepackage{tikz}
\usepackage{amsmath}
\usepackage{lmodern}
\usepackage{microtype}
\usepackage{hyperref}
\tcbuselibrary{skins, breakable, listings}

\definecolor{primary}{HTML}{1A237E}
\definecolor{accent}{HTML}{3949AB}
\definecolor{highlight}{HTML}{E8EAF6}
\definecolor{warning}{HTML}{E65100}
\definecolor{warningbg}{HTML}{FFF3E0}
\definecolor{codegray}{HTML}{F5F5F5}
\definecolor{subtitle}{HTML}{546E7A}
\definecolor{rulegray}{HTML}{B0BEC5}
\definecolor{secbg}{HTML}{E3F2FD}
\definecolor{secaccent}{HTML}{1565C0}

\definecolor{greenprimary}{HTML}{1B5E20}
\definecolor{greenaccent}{HTML}{2E7D32}
\definecolor{greenhighlight}{HTML}{E8F5E9}

\tcbset{
  categorybox/.style={
    enhanced, breakable,
    colback=greenhighlight,
    colframe=greenaccent,
    colbacktitle=greenprimary,
    coltitle=white,
    fonttitle=\bfseries\large,
    arc=6pt,
    boxrule=1.2pt,
    titlerule=0pt,
    attach boxed title to top left={yshift=-3mm, xshift=8mm},
    boxed title style={colframe=greenprimary, arc=4pt, boxrule=0.8pt},
    left=8pt, right=8pt, top=10pt, bottom=6pt,
    before skip=14pt, after skip=6pt,
  },
  warningbox/.style={
    enhanced,
    colback=warningbg,
    colframe=warning,
    arc=5pt,
    boxrule=1pt,
    left=6pt, right=6pt, top=4pt, bottom=4pt,
  },
}

\titleformat{\section}{\Large\bfseries\color{greenprimary}}
  {\thesection}{0.8em}{}[\color{rulegray}\titlerule]
\titlespacing{\section}{0pt}{18pt}{8pt}

\pagestyle{fancy}
\fancyhf{}
\fancyhead[L]{\small\color{subtitle}ООП — Въпроси за изпит}
\fancyhead[R]{\small\color{subtitle}C++}
\fancyfoot[C]{\small\color{subtitle}\thepage}
\renewcommand{\headrulewidth}{0.4pt}
\renewcommand{\headrule}{\color{rulegray}\hrule width\headwidth height\headrulewidth}

\setlist[itemize,1]{
  label={\color{greenaccent}\faChevronRight},
  leftmargin=1.4em,
  itemsep=3pt,
  parsep=0pt,
}
\setlist[itemize,2]{
  label={\color{greenaccent!70!white}\faCircle[regular]},
  leftmargin=1.8em,
  itemsep=2pt,
}

\newcommand{\nope}[1]{%
  \begin{tcolorbox}[warningbox]
    \faExclamationTriangle\quad\textbf{\color{warning}Важно:} #1
  \end{tcolorbox}%
}
\newcommand{\qtag}[1]{\texttt{\color{greenaccent}\small#1}}

\begin{document}
\selectlanguage{bulgarian}

\begin{tikzpicture}[remember picture, overlay]
  \fill[greenprimary] (current page.north west) rectangle
    ([yshift=-3.6cm]current page.north east);
  \node[anchor=west, white, font=\Huge\bfseries]
    at ([xshift=2.4cm, yshift=-1.3cm]current page.north west)
    {Въпроси по ООП — C++};
  \node[anchor=west, greenhighlight!85!greenprimary, font=\large]
    at ([xshift=2.4cm, yshift=-2.4cm]current page.north west)
    {Въпроси от изпита лятна сесия за подготовка за изпит септемврийска сесия};
\end{tikzpicture}

\vspace{1.8cm}

\begin{tcolorbox}[categorybox, title={1.\ Основни принципи на ООП въпроси}]
\begin{itemize}
  \item Абстракция — дефиниция; какво представлява; пример за абстракция
  \item Енкапсулация — дефиниция; ролята на модификаторите за достъп (\qtag{private}, \qtag{public}, \qtag{protected}) и защо са важни
  \item Наследяване — разлика между \qtag{std} и наследяване 
  \nope{Според мен не е разбран въпроса за \qtag{std} и наследяване, понеже \qtag{std} е библиотека!}
  \item Наследяване \textbf{продължение} - дефиниция; ред на извикване на конструктори и деструктори в контекста на наследяването
  \item Полиморфизъм — дефиниция; compile-time и runtime; конкретен пример като се използват шаблони, зададен от асистент (примерите са за малко по-висока оценка)
\end{itemize}
\end{tcolorbox}

\begin{tcolorbox}[categorybox, title={2.\ Класове и конструктори}]
\begin{itemize}
\item Ключова дума \qtag{class}
  \item Видове конструктори — \qtag{Default Constructor}, \qtag{Parameterized Constructor}, \qtag{Copy Constructor}, \qtag{Move Constructor}, \qtag{Conversion Constructor / Explicit}, \qtag{$=$delete Constructor}, \qtag{Implicit Default Constructor}
  \item Инициализиращ списък — какво представлява и какви са предимствата му пред това да инициализираме член-данните в scope-a
  \item Explicit конструктор — кога го използваме и защо
  \item Implicit конструктор — кога го използваме и защо
  \item Инварианти на клас — какво представляват, примери за инварианти на клас
  \item Указателят \qtag{this} — каква е неговата роля и за какво го използваме
  \nope{(За по-висока оценка) Как е имплементиран, как се възприема от компилатора!}
  \item Модификаторите за достъп (\qtag{private}, \qtag{public}, \qtag{protected}) при наследяване — кога се ползват и какво прави всеки един от тях
  \item Агрегация и Композиция — основни разлики, кой указател се ползва при всяка (raw pointer при агрегация)
  \item Pед на извикване на конструктори и деструктори
  \item Освобождаване на памет - ръчно и чрез STL контейнери, разлики
  \item \textbf{Rule of 3 / Rule of 5} — кога се имплементират, какво представляват (какво съдържат)
\end{itemize}

\nope{Виртуален конструктор — \textbf{не съществува} в C++!}
\nope{Ключова дума \texttt{interface} — \textbf{не съществува} в C++!}
\nope{Параметри в деструктор — \textbf{не можем да подаваме} и деструкторът не хвърля изключения (може да се пита защо)!}
\end{tcolorbox}

\begin{tcolorbox}[categorybox, title={3.\ Деструктори}]
\begin{itemize}
  \item Видове деструктори — \qtag{Default Destructor}, \qtag{Virtual Destructor}, \qtag{Pure Virtual Destructor}, \qtag{Deleted Destructor} 
  \item Виртуален деструктор \qtag{Virtual Destructor} — защо ни е нужен, кога се ползва
  \item Private / protected деструктор — какво се случва при опит за създаване в \qtag{main}
  \nope{%
  Сценарии: \\
  1. Създаване на обекта на стека \\
  2. Създаване на обекта в динамичната памет%
  }
  \item Как да избегнем double deletion
  \item RAII — какво е RAII, примери за RAII
  \item RVO — какво е RVO, примери за RVO
\end{itemize}
\end{tcolorbox}

\begin{tcolorbox}[categorybox, title={4.\ Виртуални функции и полиморфизъм}]
\begin{itemize}
  \item Разлика между \qtag{virtual} и \qtag{pure virtual} функции
  \item Виртуална таблица (vtable) — как изглежда, какво представлява
  \nope{Виртуална таблица — някои асистенти разписват класове и фукции и се иска да се каже как изглежда виртуалната таблица}
  \item Object slicing — какво е оbject slicing, защо е проблем оbject slicing-а
  \item Override и overload на функции — разлика между оverride и overload на функции, примери за оverride и overload на функции
  \item Абстракция и интерфейси в C++
\end{itemize}
\end{tcolorbox}

\begin{tcolorbox}[categorybox, title={5.\ Шаблони (Templates)}]
\begin{itemize}
\item Какво са Templates и какви проблеми решават и съответно какви проблеми създават
  \item Шаблонна специализация — примери и какво представлява
  \item За шаблоните асистент е питал, ако имаме някакъв оператор и направим шаблон и някакъв оператор примерно, който не е дефиниран за даден тип, ако се опитаме да го извикаме с този тип какво се случва - грешка.
  \item Написана е била шаблонна фукция на дъската и \qtag{f(string)}, \qtag{f(int)}, \qtag{f(double)} — питано е било какво се случва в този конкретен пример
  \item Concepts — какво са сoncepts и за какво служат
  \item Compile-time полиморфизъм — искало се е конкретен пример с шаблони
  \item SFINAE — какво представлява и кое е по-доброто
\end{itemize}
\end{tcolorbox}

\begin{tcolorbox}[categorybox, title={6.\ Умни указатели (Smart Pointers)}]
\begin{itemize}
  \item \qtag{unique\_ptr}
    \begin{itemize}
    \item Дефиниция и какво представлява (Ownership)
      \item Имплементация от нулата, как бихме си направили наш \qtag{unique\_ptr}
      \item Защо copy constructor и \qtag{operator$=$} са \qtag{$=$delete}
      \item Разлика между \qtag{unique\_ptr} и \qtag{shared\_ptr} относно \qtag{use\_count}
    \end{itemize}
  \item \qtag{shared\_ptr}
    \begin{itemize}
     \item Дефиниция и какво представлява (Ownership)
      \item Control block и reference count — какво са control block и reference count, къде в паметта се пазят control block и reference count
      \item Use count — как работи
    \end{itemize}
  \item \qtag{weak\_ptr}
    \begin{itemize}
    \item Дефиниция и какво представлява (Ownership)
      \item За какво се ползва
      \item Разлики от \qtag{shared\_ptr}
    \end{itemize}
  \item Ownership — какво означава и как ни помага
  \item Raw pointer при агрегация — кога и защо се използва
\end{itemize}
\end{tcolorbox}

\begin{tcolorbox}[categorybox, title={ 7.\ Ключови думи}]
\begin{multicols}{2}
\begin{itemize}
\item Ключова дума \qtag{class}
  \item Ключова дума \qtag{const} в началото и в края на функция
  \item Ключова дума \qtag{const} променлива
  \item Ключова дума \qtag{static} — член-данни и функции
  \item Ключова дума \qtag{mutable} - какво представлява
  \item Ключова дума \qtag{auto} - какво представлява
  \item Ключова дума \qtag{decltype} — разлика с \qtag{auto}
  \item Ключова дума \qtag{friend} - какво представлява
  \item Ключова дума \qtag{explicit} - какво представлява
\end{itemize}
\end{multicols}
\nope{Ключова дума \qtag{auto} - имало е въпрос свързан с това, кога се разбира какъв реално е типа!}
\end{tcolorbox}

\begin{tcolorbox}[categorybox, title={8.\ Move семантика и стойности}]
\begin{itemize}
  \item lvalue и rvalue — разлика между двете
  \item rvalue reference (\qtag{\&\&}) - какво представлява
  \item \qtag{std::move} и move semantics - какво представляват и кога се използват
  \item \qtag{std::forward} (perfect forwarding)
  \item Живот на обект — какво представлява понятието
\end{itemize}
\end{tcolorbox}

\begin{tcolorbox}[categorybox, title={9.\ Изключения}]
\begin{itemize}
\item Какво са exceptions и защо се използват?
\item Какво e stack unwinding?
  \item Видове exceptions
  \item Нива на exception safety (\qtag{No-throw}, \qtag{Strong}, \qtag{Basic}, \qtag{No guarantee})
  \item Кoe e по-добро от exceptions и защо? (Hint: \qtag{std::expected})
\end{itemize}
\end{tcolorbox}

\begin{tcolorbox}[categorybox, title={10.\ STL типове и собствени имплементации}]
\begin{multicols}{2}
\begin{itemize}
  \item \qtag{std::string} и \qtag{char*} - кое е по-добро, защо и как бихме си имплементирали myString?
  \nope{За собствената имплементация: Трябва класът да има char*, size, capacity, rule of 5, resize функция!}
  \item \qtag{std::vector} и \qtag{std::list} — кога кое ще използваме? Предимства и недостатъци
  \item \qtag{std::optional} - Какво представлява? Предимства и недостатъци
  \item \qtag{std::expected} - Какво представлява? Предимства и недостатъци
  \item \qtag{std::variant} - Какво представлява? Предимства и недостатъци
  \item \qtag{std::tuple} - Какво представлява? Предимства и недостатъци
\end{itemize}
\end{multicols}
\end{tcolorbox}

\begin{tcolorbox}[categorybox, title={11.\ Дизайн принципи и шаблони}]
\begin{itemize}
  \item \textbf{SOLID} — всичките пет принципа; 
  \item \qtag{S} от \qtag{SOLID} (Single Responsibility) по-подробно
  \item \textbf{Builder pattern} — какво представлява, примери
  \item \textbf{Prototype design pattern} — какво представлява, примери
  \item \textbf{Singleton} — какво представлява, примери, как се имплементира, какви ограничения налага
\end{itemize}
\end{tcolorbox}

\begin{tcolorbox}[categorybox, title={12.\ Cast, Namespace, Enum, Lambda}]
\begin{itemize}
  \item \qtag{dynamic\_cast} и \qtag{static\_cast} — основни разлики и употреба 
  \item \qtag{namespace} — какво представлява, с какво е полезно
  \item \qtag{enum} и \qtag{enum class} — основни разлики и употреба
  \item Lambda функции — какво представлява, какво означават \texttt{[=]}, \texttt{[\&]}, подаване по стойност в Lambda функции
\end{itemize}
\end{tcolorbox}

\begin{tcolorbox}[categorybox, title={13.\ Памет, компилация и свързване}]
\begin{itemize}
  \item Heap и Stack — разлики в двата вида памет
  \item Разделна компилация — да се обясни целият процес (preprocessor → compiler → linker)
  \item Задача с linker грешка — да се разпознае и да се даде обяснение
  \item Предефиниране на оператори (\textit{operator overloading})
  \item оператор spaceship \qtag{$<=>$} - защо го използваме, какво спестява и какво означава да е $=$default
\end{itemize}
\end{tcolorbox}

\begin{tcolorbox}[categorybox, title={14.\ Практически задачи}]
\begin{itemize}
  \item Имплементирайте \qtag{MyString}: \qtag{char*}, \qtag{size}, \qtag{capacity}, Rule of 5, \qtag{resize()}
  \item Имплементирайте \qtag{unique\_ptr}
  \item Нарисувайте vtable за дадена йерархия от класове
  \item Дадена е шаблонна функция + \qtag{f(string)}/\qtag{f(int)}/\qtag{f(double)} — какво ще се извиква
  \item Задача с linker грешка — намерете и обяснете грешката
\end{itemize}
\end{tcolorbox}

\vfill
\begin{center}
  \color{subtitle}\small
  \faInfoCircle\quad Всички въпроси са събрани от реални изпити. Мога да изпратя и разписана теория от лекции. Акцентирайте върху \textbf{практическите примери}.
\end{center}

\end{document}
